\section{First-order coefficient currents below the quadratic energy threshold}
\label{sec:r11-first-order-power-currents}

Retain the original arbitrary background, affinoid domain,
finite power-energy family, and notation (PG1)--(PG2).
The construction requires only $\sigma>0$; it does not require
that every quadratic energy current exist.
Set $k=|\alpha|\leq q-1$ and $\kappa=q+\sigma$.
A first-order \emph{coefficient measure} is its pairing with
each fixed original complementary smooth form, with actual
Borel scalar multipliers inserted. All variation bounds
are on compact original test stages.

\begin{lemma}[Bounded first-order chain and Borel locality]
\label{lem:r11-first-order-borel-locality}
Let $Q$ be an arbitrary original positive separately closed
current and $v_1,\ldots,v_m$ bounded classical psh functions.
For every $C^1$ function $F$ near their bounded range box,
\[
 d'(F(\mathbf v)Q)=\sum_i F_i(\mathbf v)d'(v_iQ),\qquad
 d''(F(\mathbf v)Q)=\sum_i F_i(\mathbf v)d''(v_iQ).
 \tag{FO1}
\]
These are coefficient-Radon identities. For any real $t$,
\[
 d'(\max(v_i,t)Q)=\mathbf1_{\{v_i>t\}}d'(v_iQ),\qquad
 d''(\max(v_i,t)Q)=\mathbf1_{\{v_i>t\}}d''(v_iQ).
 \tag{FO2}
\]
If positive separately closed backgrounds $Q,Q'$ of the
same degree have equal coefficient restrictions on a Borel
set $A$ in the full original test sense, then for every
fixed bounded classical $v$,
\[
 \mathbf1_A d'(vQ)=\mathbf1_A d'(vQ'),\qquad
 \mathbf1_A d''(vQ)=\mathbf1_A d''(vQ').
 \tag{FO3}
\]
Consequently $d'(u_{i,j}T_{\alpha,j})$ and
$d''(u_{i,j}T_{\alpha,j})$ are consistent on $E_j$
as $j$ increases.
\end{lemma}

\noindent\textbf{Status: \texttt{DIRECT\_COROLLARY}} of
the bounded first-order measure theorem and strict maximum
locality, without differentiating a Borel restriction.

\begin{proof}
Take a common local bounded smooth decreasing approximation
of the tuple. Lemma~\ref{lem:r11-bounded-first-order}
gives local coefficient total-variation convergence of each
first derivative. The bounded continuous multipliers
$F_i(\mathbf v_l)$ converge pointwise. Dominated convergence
against the fixed limiting derivative measures, together
with the variation errors, passes the smooth chain formula.
Bounded convergence against the coefficients of $Q$
identifies its left side. This proves (FO1).

Choose smooth convex nondecreasing hinge approximations
$H_\epsilon$, with uniform error at most $\epsilon$,
$0\leq H_\epsilon'\leq1$, and
$H_\epsilon'(s)\to\mathbf1_{\{s>t\}}$ at every $s$,
including value zero at $t$. Formula (FO1) and dominated
convergence against the fixed derivative measure prove (FO2).

For (FO3), take one common local smooth approximation $v_l$
independent of the backgrounds. Closedness gives
$d'(v_lQ)=d'v_l\wedge Q$, whose restriction on $A$
equals that for $Q'$ in every actual refinement containing
the approximation's coordinates. Pass this equality through
coefficient total-variation convergence.
For consistency on $E_j$, first change $T_{\alpha,l}$
to $T_{\alpha,j}$ by bounded maximum locality and (FO3),
holding $u_{i,l}$ fixed. Then apply (FO2) to
$u_{i,j}=\max(u_{i,l},-j)$. The double-prime proofs are identical.
\end{proof}

\begin{lemma}[A first-order capacity bound]
\label{lem:r11-first-order-capacity-bound}
Let $-L\leq b_i\leq0$ be bounded classical psh functions,
and let $Q_\alpha$ be their bounded mixed product over $S$,
$|\alpha|=k\leq q-1$. For every original compact
complementary first-order test $\eta$ and Borel $A\subset V$,
\[
 \bigl|\mu_{d'(b_iQ_\alpha),\eta}\bigr|(A)
                  \leq C_\eta L^{k+1}C_S(A;V).
 \tag{FO4}
\]
There is the analogous double-prime estimate. The constant
is independent of the bounded potentials and $L>0$.
\end{lemma}

\noindent\textbf{Status: \texttt{NEW\_PROVED}}.

\begin{proof}
Put $d=q-k$ and choose a compact cutoff and fixed original
moment control $\beta$ for $\eta$. The restricted (NP1),
after taking a supremum over Borel multipliers of absolute
value at most one, gives
\[
 \bigl|\mu_{d'(b_iQ_\alpha),\eta}\bigr|(A)^2
 \leq C_\eta
       \left(\int_A\chi\,\Gamma_{Q_\alpha}(b_i)
                                  \wedge\beta^{d-1}\right)
       \left(\int_A\chi\,\beta^d\wedge Q_\alpha\right).
 \tag{FO5}
\]
The factors are bounded by $CL^{k+2}C_S(A;V)$ from
(PG8) and by $C'L^kC_S(A;V)$ from (DG3).
Their square root proves (FO4). The coefficient controls
remain valid on every actual refinement.
\end{proof}

\begin{theorem}[Intrinsic unbounded first-order coefficient currents]
\label{thm:r11-unbounded-first-order}
For $\sigma>0$, $|\alpha|=k\leq q-1$, and every index $i$,
the intrinsic distributional currents
\[
 {\cal J}'_{\alpha,i}=d'(u_iN_\alpha),\qquad
 {\cal J}''_{\alpha,i}=d''(u_iN_\alpha)
 \tag{FO6}
\]
are coefficient-Radon. They have zero coefficient variation
on the pole union, and
\[
 d'(u_{i,j}T_{\alpha,j})\longrightarrow{\cal J}'_{\alpha,i},
 \qquad
 d''(u_{i,j}T_{\alpha,j})\longrightarrow{\cal J}''_{\alpha,i}
 \tag{FO7}
\]
in local coefficient total variation.
For every original compact complementary test $\eta$,
$a\geq0$ with $a+k+1\leq\kappa$, and $R\geq1$,
\[
 \begin{split}
 \sup_j\int_{\{w\leq-R\}}W_j^a
       \bigl|d\mu_{d'(u_{i,j}T_{\alpha,j}),\eta}\bigr|
     &\leq C_\eta R^{a+k+1-\kappa}{\cal J}_w(R/4),\\
 \int_{\{w\leq-R\}}W^a
       \bigl|d\mu_{{\cal J}'_{\alpha,i},\eta}\bigr|
     &\leq C_\eta R^{a+k+1-\kappa}{\cal J}_w(R/4).
 \end{split}
 \tag{FO8}
\]
Full local weighted variations are finite, the double-prime
bounds are identical, and a continuous multiplier of growth
at most $W^a$ may be inserted in (FO7).
These assertions apply even when $k+2>q+\sigma$, where
the graded theorem constructs no quadratic energy current.
\end{theorem}

\noindent\textbf{Status: \texttt{NEW\_PROVED}}.
These are first-order coefficient measures, not closed
positive backgrounds for a subsequent BT iteration.

\begin{proof}
On $A_T=\{-2T<w\leq-T\}$ further truncate every current
slot at $-2T$. Lemma
\ref{lem:r11-first-order-borel-locality} permits this in
both the derivative slot and its closed background.
Apply (FO4) with $L=2T$. Since $W_j\leq1+2T$ on the shell,
\[
 \int_{A_T}W_j^a
       \bigl|d\mu_{d'(u_{i,j}T_{\alpha,j}),\eta}\bigr|
                    \leq C T^{a+k+1}F_w(T/2).
 \tag{FO9}
\]
Every finite-stage derivative is zero on the pole union
by (FO4). Summing (FO9) with (PG6) proves the first (FO8).

The cut coefficient measures are consistent on $E_j$.
Their variations increase and are locally bounded by
the preceding estimates. Hence they define Radon limits
on the finite-value locus, extended by zero on the
pole union. The second (FO8) holds for these limits.
Agreement on $E_j$ and the two tails at $R=j$ prove
coefficient total-variation convergence of the uncut
derivatives and all the stated weighted versions.
Finite linear coefficient identities and actual
refinement identities pass to these limits.

Equation (PG4) with $F(x)=x_i$ gives
$u_{i,j}T_{\alpha,j}\to u_iN_\alpha$ on every original
test. Distributional continuity identifies the limits
with (FO6), including their current continuity on
each original compact stage. The double-prime proof
is identical.
\end{proof}

\begin{theorem}[Unbounded first-order chain and conormal identities]
\label{thm:r11-unbounded-first-order-chain}
Let $F$ be $C^1$ near the nonpositive orthant and suppose
that for $1\leq r\leq\kappa-k$,
\[
 |F(x)|\leq A W(x)^r,\qquad
 |F_i(x)|\leq A W(x)^{r-1}.
 \tag{FO10}
\]
Then, on all original complementary tests,
\[
 d'(F(\mathbf u)N_\alpha)
                  =\sum_iF_i(\mathbf u){\cal J}'_{\alpha,i},
 \qquad
 d''(F(\mathbf u)N_\alpha)
                  =\sum_iF_i(\mathbf u){\cal J}''_{\alpha,i}.
 \tag{FO11}
\]
If $F(\mathbf u)$ equals a constant on the actual
finite-value locus of an ordinary germ, then
\[
       \sum_iF_i(\mathbf u){\cal J}'_{\alpha,i}
       =\sum_iF_i(\mathbf u){\cal J}''_{\alpha,i}=0
 \tag{FO12}
\]
on that germ and every actual moment refinement.
\end{theorem}

\noindent\textbf{Status: \texttt{NEW\_PROVED}} for
(FO11); (FO12) is a \texttt{DIRECT\_COROLLARY}.
The conormal equation uses an actual $C^1$ value
relation, not an analytic delta-form relation quotient.

\begin{proof}
Apply (FO1) to common bounded truncations. The multiplier
on the left has total exponent $r+k\leq\kappa$, so
(PG4) passes it to its coefficient-defined limit.
On the right the weight has growth $W^{r-1}$ and the
derivative budget is $k+1$, giving the same exponent.
Thus (FO8) passes each term in coefficient variation;
distributional continuity gives (FO11).
If the actual value relation is $F(\mathbf u)=c$,
then $F(\mathbf u)N_\alpha=cN_\alpha$, because
$N_\alpha$ is zero on the omitted poles.
Separate closedness now gives (FO12).
\end{proof}

\paragraph{Relation to the analytic relation quotient.}
This supplies actual coefficient derivatives and conormal
identities for every original test and refinement with
arbitrary original $S$. It does not supply an arbitrary
wedge of these measures, closedness of level restrictions,
or equality on additional GK delta tests. The complete
signed PS/PL quotient and the low-real-degree constant-tensor
order gate remain \texttt{OPEN}.
